\documentclass[12pt,a4paper]{article}

% ---- Encoding and language ----
\usepackage[utf8]{inputenc}
\usepackage[T1]{fontenc}
\usepackage[brazil]{babel}

% ---- Page geometry (ABNT margins) ----
\usepackage[top=3cm, bottom=2cm, left=3cm, right=2cm]{geometry}

% ---- Typography ----
\usepackage{setspace}
\onehalfspacing
\usepackage{parskip}
\setlength{\parindent}{1.25cm}
\setlength{\parskip}{0pt}

% ---- Math ----
\usepackage{amsmath}
\usepackage{amssymb}

% ---- Tables ----
\usepackage{booktabs}
\usepackage{array}

% ---- Code listings ----
\usepackage{listings}
\usepackage{xcolor}
\lstset{
    basicstyle=\ttfamily\small,
    keywordstyle=\color{blue}\bfseries,
    commentstyle=\color{gray}\itshape,
    stringstyle=\color{teal},
    numbers=left,
    numberstyle=\tiny\color{gray},
    numbersep=8pt,
    frame=single,
    breaklines=true,
    tabsize=4,
    showstringspaces=false,
    captionpos=b
}

% ---- Graphics ----
\usepackage{graphicx}

% ---- Hyperlinks (load last) ----
\usepackage[hidelinks]{hyperref}

% ---- Bibliography ----
% Using traditional BibTeX (compatible with latexmk -bibtex)
\bibliographystyle{plain}

% ========================================================
\begin{document}

% ---- Cover page ----
\begin{titlepage}
    \centering
    \vspace*{2cm}

    {\large Universidade Federal de Santa Catarina}\\[0.3cm]
    {\large Departamento de Informática e Estatística}\\[0.3cm]
    {\large Modelagem e Simulação}

    \vspace{3cm}

    {\LARGE \textbf{Desenvolvimento de Componente de Análise de Dados no GenESyS}}\\[0.5cm]
    {\large Tema 11.3.1 -- Foco em Diagnósticos Adicionais e Ajustes de Curvas}

    \vspace{3cm}

    {\large \textbf{Autor:} Gabriel Salmoria}

    \vfill

    {\large Florianópolis, \the\year}
\end{titlepage}

% ---- Table of contents ----
\tableofcontents
\newpage

% ---- Chapters ----
\section{Introdução}

\subsection{Contextualização}

O GenESyS (\textit{Generic Extensible Simulator}) é uma plataforma de modelagem e simulação de eventos discretos desenvolvida em C++, cujo código-fonte é mantido em repositório público no GitHub \cite{genesys-github}. A plataforma adota uma arquitetura extensível baseada em interfaces abstratas e implementações concretas intercambiáveis, permitindo a adição de novos componentes sem alteração do núcleo do simulador.

No contexto de simulação, a análise estatística dos resultados produzidos por modelos é uma etapa fundamental do processo científico. Após a execução de um modelo, é necessário verificar a distribuição dos dados de saída, ajustar distribuições teóricas, aplicar testes de aderência e realizar inferências sobre as populações simuladas \cite{law2015}. Até o início deste trabalho, o pacote \texttt{source/tools/} do GenESyS já oferecia implementações para ajuste de distribuições (\texttt{Fitter\_if}), testes de hipóteses (\texttt{Hpothesis\-Tester\_if}) e resolução numérica (\texttt{Solver\_if}), porém sem uma fachada unificada que integrasse todas essas capacidades em um único ponto de acesso.

\subsection{Objetivo}

O objetivo deste trabalho é implementar o componente \texttt{DataAnalyzer}, uma ferramenta de análise de dados integrada ao GenESyS que oferece:

\begin{itemize}
    \item carregamento de dados a partir de memória ou de arquivos texto nos formatos suportados pelo simulador;
    \item estatísticas descritivas completas (média, mediana, moda, quartis, assimetria, curtose, entre outras);
    \item estruturas exploratórias de histograma e boxplot;
    \item ajuste de sete distribuições de probabilidade com seleção automática da melhor;
    \item testes de aderência qui-quadrado e Kolmogorov-Smirnov;
    \item análise de séries temporais (média móvel e autocorrelação);
    \item intervalos de confiança e testes de hipóteses para uma e duas populações.
\end{itemize}

O tema enquadra-se na diretriz 11.3.1 da proposta -- \textit{foco maior em diagnósticos adicionais e ajustes de curvas} -- e é referenciado no GenESyS como \texttt{DataAnalyzer}.

\subsection{Justificativa}

A ausência de uma interface unificada de análise forçava o usuário a instanciar e coordenar manualmente \texttt{Fitter\_if}, \texttt{Hypothesis\-Tester\_if} e parsers de arquivo, além de reimplementar cálculos estatísticos básicos em cada contexto de uso. Um componente façade reduz o acoplamento entre camadas do simulador e serve como base para futuros painéis gráficos de análise de resultados.

Do ponto de vista pedagógico, o trabalho exigiu compreensão profunda da arquitetura do GenESyS -- especialmente o padrão \texttt{TraitsTools<T>} de injeção de implementações -- e a integração de métodos numéricos já existentes na plataforma (integração por quadratura, busca de raízes por bissecção) para suportar novas funcionalidades.

\newpage

\section{Desenvolvimento}

\subsection{Visão geral da arquitetura}

O componente foi desenvolvido em três camadas, seguindo o padrão arquitetural já adotado pelo GenESyS:

\begin{enumerate}
    \item \textbf{Interface abstrata} (\texttt{DataAnalyzer\_if.h}) -- define os tipos de saída e os métodos puros virtuais;
    \item \textbf{Implementação concreta} (\texttt{DataAnalyzerDefaultImpl1}) -- realiza todos os cálculos e delega para \texttt{FitterDefaultImpl} e \texttt{HypothesisTesterDefaultImpl1};
    \item \textbf{Registro de traits} (\texttt{TraitsTools<DataAnalyzer\_if>}) -- vincula a interface à implementação concreta, tornando a substituição transparente.
\end{enumerate}

A Tabela~\ref{tab:arquivos} lista os arquivos criados ou modificados.

\begin{table}[htbp]
\centering
\caption{Arquivos modificados ou criados}
\label{tab:arquivos}
\begin{tabular}{ll}
\hline
\textbf{Arquivo} & \textbf{Ação} \\
\hline
\texttt{source/tools/Fitter\_if.h}              & Adicionado \texttt{setDataValues()} \\
\texttt{source/tools/FitterDefaultImpl.h}        & Adicionados \texttt{setDataValues()} e \texttt{\_computeStatsFromLoadedData()} \\
\texttt{source/tools/FitterDummyImpl.h/.cpp}     & Adicionado stub vazio \\
\texttt{source/tools/DataAnalyzer\_if.h}         & Novo (interface + structs de saída) \\
\texttt{source/tools/DataAnalyzerDefaultImpl1.h} & Novo (declaração da classe) \\
\texttt{source/tools/DataAnalyzerDefaultImpl1.cpp} & Novo ($\approx$810 linhas, toda a lógica) \\
\texttt{source/tools/TraitsTools.h}              & Adicionada especialização para \texttt{DataAnalyzer\_if} \\
\texttt{source/tools/CMakeLists.txt}             & Adicionado \texttt{DataAnalyzerDefaultImpl1.cpp} \\
\texttt{source/tools/main.cpp}                   & Adicionado comando \texttt{-{}-demo} \\
\hline
\end{tabular}
\end{table}

\subsection{Extensão da interface Fitter\_if (M1)}

\texttt{FitterDefaultImpl} só aceitava dados via arquivo binário (\texttt{setDataFilename}). Para que o \texttt{DataAnalyzer} pudesse repassar ao fitter um vetor já carregado em memória, foi adicionado o método \texttt{setDataValues(const std::vector<double>\&)} à interface \texttt{Fitter\_if} e à implementação concreta.

A lógica interna de \texttt{FitterDefaultImpl} foi refatorada: o código que computava estatísticas de cache foi extraído para o método privado \texttt{\_computeStatsFromLoadedData()}, chamado tanto pelo caminho original de arquivo binário quanto pelo novo caminho de memória. Os flags de cache \texttt{\_cacheLoaded} e \texttt{\_cacheUsable} são definidos por ambos os caminhos, de modo que o mecanismo de \textit{lazy loading} existente em \texttt{\_ensureDataLoaded()} enxerga os dados como já carregados e não tenta abrir arquivo algum.

\subsection{Interface DataAnalyzer\_if e structs de saída (M2)}

A interface \texttt{DataAnalyzer\_if} define cinco structs de saída puramente de dados (sem métodos):

\begin{itemize}
    \item \texttt{SummaryStatistics} -- n, mín, máx, amplitude, média, mediana, moda, Q1, Q3, variância, desvio-padrão, CV, assimetria, curtose;
    \item \texttt{HistogramData} -- número de classes, limites inferiores, frequências absolutas e relativas;
    \item \texttt{BoxplotData} -- mín, Q1, mediana, Q3, máx, IIQ, lista de \textit{outliers};
    \item \texttt{FitResult} -- nome da distribuição, SSE, até cinco parâmetros, flag de validade;
    \item \texttt{GoFResult} -- nome, estatística de teste, p-valor, valor crítico, nível de significância, decisão, conclusão textual.
\end{itemize}

A grafia americana (\textit{Analyzer}) foi adotada para evitar confusão com o antigo \texttt{DataAnalyser\_if} (grafia britânica), que possui dependência de \texttt{ExperimentManager\_if} e pertence a outra camada arquitetural.

\subsection{Implementação: carregamento de dados e estatísticas descritivas (M3)}

O carregamento de dados suporta dois caminhos:

\begin{itemize}
    \item \texttt{setDataValues(vector<double>)}: armazena os dados em memória, ordena o vetor em \texttt{\_sortedData} e computa média, variância e desvio-padrão em uma única passagem (\texttt{\_rebuildCache()});
    \item \texttt{setDataFilename(string)}: utiliza \texttt{SimulationResultsDatasetParser::loadFromTextFile} para suportar os formatos de texto do GenESyS (RawNumeric, RecordLegacy, RecordEnriched, GuiTabular), depois chama \texttt{setDataValues}.
\end{itemize}

Os quantis são calculados pelo método R7 (interpolação linear): dado percentil $p$, calcula-se $h = (n-1)p$, $i = \lfloor h \rfloor$, $f = h - i$, e o resultado é $x_{(i)}(1-f) + x_{(i+1)}f$. Assimetria e curtose são calculadas diretamente como momentos padronizados de terceira e quarta ordem (curtose em excesso, subtraindo 3).

\textbf{Decisão de projeto -- \texttt{StatisticsDataFile\_if}:} A proposta sugeria encaminhar as estatísticas descritivas pela interface \texttt{StatisticsDatafile\_if}. Essa interface estende \texttt{CollectorDatafile\_if}, que é estruturalmente acoplada a um arquivo binário de coletor -- não há construtor nem setter que aceite \texttt{vector<double>}. Utilizá-la exigiria escrever os dados em arquivo binário e relê-los, o que derrota o propósito de operação em memória. A decisão foi calcular as estatísticas diretamente, produzindo resultados equivalentes sem o custo de I/O.

\subsection{Ajuste de distribuições (M4)}

O método \texttt{fitDistribution(name)} despacha para o método correto de \texttt{FitterDefaultImpl} com base no nome da distribuição e empacota SSE e parâmetros em \texttt{FitResult}. As sete distribuições suportadas e seus parâmetros são:

\begin{center}
\texttt{uniform} (mín, máx), \texttt{triangular} (mín, moda, máx), \texttt{normal} (média, dp), \texttt{exponential} (média), \texttt{erlang} (média, m), \texttt{beta} ($\alpha$, $\beta$, inf, sup), \texttt{weibull} (forma, escala).
\end{center}

O método \texttt{fitAll()} chama \texttt{\_fitter.fitAll()} para obter o nome da melhor distribuição e, em seguida, chama \texttt{fitDistribution(nome)} para recuperar os parâmetros. O duplo ajuste é intencional: \texttt{fitAll} retorna apenas nome e SSE.

\subsection{Testes de aderência (M5 e M6)}

\textbf{Qui-quadrado (M5):} O número de classes é determinado pela fórmula de Sturges ($k = \max(5,\ 1 + \lceil \log_2 n \rceil)$). Classes com frequência esperada $E < 5$ são fundidas com a adjacente. O grau de liberdade é $k_{\text{fundido}} - 1 - p$, onde $p$ é o número de parâmetros estimados a partir dos dados. O p-valor é calculado integrando numericamente a PDF da distribuição qui-quadrado usando \texttt{SolverDefaultImpl1}; o valor crítico é obtido por bissecção sobre a CDF acumulada.

\textbf{Kolmogorov-Smirnov (M6):} A estatística $D_n = \max_i |F_n(x_i) - F(x_i)|$ é calculada sobre as estatísticas de ordem. O p-valor usa a série assintótica de Kolmogorov com 200 termos ($2\sum_{k=1}^{200}(-1)^{k+1}e^{-2k^2(\sqrt{n}\,D_n)^2}$); o valor crítico usa a fórmula aproximada $D_\alpha \approx \sqrt{-\ln(\alpha/2)/(2n)}$. A aproximação assintótica subestima levemente o p-valor para $n < 35$.

\subsection{Análise de séries temporais (M7)}

\begin{itemize}
    \item \texttt{movingAverage(w)}: janela deslizante com soma acumulada $O(n)$; saída com $n - w + 1$ elementos;
    \item \texttt{autocorrelation(maxLag)}: estimador enviesado $\hat{r}(k) = \hat{c}(k)/\hat{c}(0)$, onde $\hat{c}(k) = \frac{1}{n}\sum_{i=0}^{n-k-1}(x_i - \bar{x})(x_{i+k} - \bar{x})$; \texttt{maxLag} é silenciosamente limitado a $n-1$;
    \item \texttt{correlogram(maxLag)}: retorna o mesmo cálculo de \texttt{autocorrelation}, diferenciado conceitualmente para uso como gráfico de barras por defasagem.
\end{itemize}

\subsection{Inferência paramétrica (M8)}

Os métodos de intervalo de confiança e teste de hipóteses delegam para \texttt{HypothesisTesterDefaultImpl1}, usando as estatísticas em cache (\texttt{\_sampleMean}, \texttt{\_sampleVariance}, \texttt{\_sampleStddev}, \texttt{\_count}, \texttt{\_confidenceLevel}). Os testes para duas populações requerem que \texttt{loadSecondSample()} tenha sido chamado previamente.

\textbf{Correção de bug:} A versão inicial de \texttt{testProportionTwoSamples} e \texttt{testVarianceTwoSamples} invocava, por engano, as sobrecargas de uma população de \texttt{HypothesisTesterDefaultImpl1}, tratando a estatística da segunda amostra como valor da hipótese nula. A correção foi passar \texttt{\_count2} como quarto argumento, selecionando as sobrecargas corretas de duas populações já existentes na implementação.

\subsection{Integração com o GenESyS}

O componente integra-se ao simulador por meio do mecanismo \texttt{TraitsTools<T>}:

\begin{lstlisting}[language=C++, caption={Registro da implementação em TraitsTools.h}]
template <> struct TraitsTools<DataAnalyzer_if> {
    typedef DataAnalyzerDefaultImpl1 Implementation;
};
\end{lstlisting}

Qualquer código do simulador pode obter uma instância concreta por meio de \texttt{TraitsTools<DataAnalyzer\_if>::Implementation}, sem depender da classe concreta diretamente. A ferramenta de linha de comando (\texttt{genesys\_tools\_application}) foi estendida com o comando \texttt{-{}-demo} para demonstrar e validar o componente sem interface gráfica.

\newpage

\section{Validação}

\subsection{Estratégia de teste}

A validação foi realizada por meio de um \textit{harness} de testes automatizados integrado ao executável \texttt{genesys\_tools\_application}, acionado pelo comando \texttt{-{}-demo}. O harness instancia \texttt{DataAnalyzerDefaultImpl1} diretamente e executa 26 verificações organizadas em 7 seções, imprimindo \texttt{[PASS]} ou \texttt{[FAIL]} para cada uma.

Três conjuntos de dados fixos são usados:
\begin{itemize}
    \item \texttt{kNormalData}: 50 valores amostrados de $N(5, 1)$;
    \item \texttt{kExpoData}: 20 valores amostrados de $\text{Exp}(2)$ (média $= 2$);
    \item \texttt{kTemporalData}: 20 valores com tendência crescente, de 1,1 a 3,0.
\end{itemize}

\subsection{Seção 1 -- Carregamento de dados}

\begin{itemize}
    \item \texttt{setDataValues} aceita 50 valores e \texttt{summaryStatistics().n == 50};
    \item \texttt{setDataFilename} carrega os mesmos 50 valores de arquivo RawNumeric temporário;
    \item \texttt{clearData} esvazie o conjunto (\texttt{n == 0});
    \item \texttt{loadSecondSample} armazena um segundo vetor de 20 valores.
\end{itemize}

\subsection{Seção 2 -- Estatísticas descritivas}

\begin{itemize}
    \item Média dentro do intervalo $[4{,}5;\ 5{,}5]$ para dados $N(5,1)$;
    \item Ordenamento: $\text{mín} < Q_1 < \text{mediana} < Q_3 < \text{máx}$;
    \item Desvio-padrão $> 0$;
    \item Consistência entre \texttt{quartile(2)}, \texttt{decile(5)} e \texttt{centile(50)} (todos devem ser iguais à mediana).
\end{itemize}

\subsection{Seção 3 -- Histograma e boxplot}

\begin{itemize}
    \item Soma das frequências absolutas igual a $n = 50$;
    \item Histograma com exatamente 8 classes;
    \item Soma das frequências relativas igual a 1,0 (tolerância $10^{-10}$);
    \item $\text{IIQ} = Q_3 - Q_1$ (verificação da fórmula do boxplot).
\end{itemize}

\subsection{Seção 4 -- Ajuste de distribuições}

\begin{itemize}
    \item Ajuste normal válido (\texttt{FitResult.valid == true});
    \item Parâmetro de média do ajuste normal dentro de $[4{,}5;\ 5{,}5]$;
    \item \texttt{fitAll} retorna resultado válido;
    \item SSE do \texttt{fitAll} $\leq$ SSE do ajuste normal (é a melhor entre todas);
    \item Média do ajuste exponencial nos dados \texttt{kExpoData} dentro de $[1{,}0;\ 3{,}5]$ (esperada $\approx 2$).
\end{itemize}

\subsection{Seção 5 -- Testes de aderência}

\begin{itemize}
    \item Teste qui-quadrado não rejeita ajuste normal a $\alpha = 0{,}05$ (dados são normais);
    \item Teste KS não rejeita ajuste normal a $\alpha = 0{,}05$.
\end{itemize}

\subsection{Seção 6 -- Séries temporais}

\begin{itemize}
    \item \texttt{movingAverage(3)} produz $n - w + 1 = 18$ valores;
    \item Primeiro valor da média móvel igual à média dos três primeiros elementos ($(1{,}1 + 1{,}3 + 1{,}2)/3$);
    \item \texttt{autocorrelation(5)} retorna \texttt{maxLag + 1 = 6} valores;
    \item $\hat{r}(0) = 1{,}0$ por construção;
    \item $\hat{r}(1) > 0{,}5$ para dados com tendência crescente;
    \item \texttt{correlogram(5) == autocorrelation(5)} (implementação compartilhada).
\end{itemize}

\subsection{Seção 7 -- Inferência paramétrica}

\begin{itemize}
    \item IC 95\% para a média contém o verdadeiro valor 5,0;
    \item Teste t bilateral para $\mu_0 = 5{,}2$ não rejeita ($\mu_0$ próximo à média amostral de 5,206);
    \item Teste t bilateral para $\mu_0 = 10{,}0$ rejeita (valor muito afastado);
    \item Limite superior do IC de variância $>$ limite inferior;
    \item Teste de médias para duas populações (normal vs.\ exponencial) rejeita $H_0: \mu_1 = \mu_2$.
\end{itemize}

\subsection{Resultado da execução}

A saída do comando \texttt{-{}-demo} é apresentada abaixo (trecho):

\begin{verbatim}
============================================================
DataAnalyzer Demo — GenESyS tools
============================================================
...
1. DATA LOADING
------------------------------------------------------------
  [PASS] setDataValues accepted 50 values
  [PASS] setDataFilename loaded same 50 values from RawNumeric text file
  [PASS] clearData empties the dataset
  [PASS] loadSecondSample stores a second vector
...
============================================================
Demo complete.
============================================================
\end{verbatim}

Todos os 26 testes passam com a implementação atual.

\subsection{Problemas encontrados e corrigidos}

Durante o desenvolvimento foram encontrados e corrigidos três problemas:

\begin{enumerate}
    \item \textbf{Incompatibilidade de tipos no CDF de Erlang:} \texttt{std::llround} retorna \texttt{long long int}, enquanto a literal \texttt{1L} é \texttt{long int}. O compilador recusou \texttt{std::max(1L, std::llround(p2))} por tipos conflitantes. Corrigido com \texttt{std::max(1LL, std::llround(p2))}.

    \item \textbf{Métodos não-const de \texttt{ConfidenceInterval}:} Os métodos \texttt{inferiorLimit()} e \texttt{superiorLimit()} de \texttt{HypothesisTester\_if::ConfidenceInterval} não são declarados \texttt{const}, portanto variáveis locais \texttt{const auto ci} falham com erro de qualificador. Corrigido declarando as variáveis CI como \texttt{auto ci} (não-const).

    \item \textbf{Testes de duas populações usando sobrecargas de uma população:} As versões iniciais de \texttt{testProportionTwoSamples} e \texttt{testVarianceTwoSamples} omitiam o argumento \texttt{\_count2}, ativando as sobrecargas de uma população e tratando a estatística da segunda amostra como valor nulo. Corrigido passando \texttt{\_count2} como quarto argumento, selecionando as sobrecargas de duas populações corretas.
\end{enumerate}

\newpage

\section{Conclusão}

O componente \texttt{DataAnalyzer} foi implementado e integrado ao GenESyS, atendendo à diretriz 11.3.1 da proposta. Partindo da necessidade de coordenar manualmente \texttt{Fitter\_if}, \texttt{HypothesisTester\_if} e parsers de arquivo, o grupo construiu uma fachada única que cobre análise descritiva, ajuste de sete distribuições com seleção automática, dois testes de aderência (qui-quadrado e KS), média móvel e autocorrelação, e inferência paramétrica para uma e duas populações.

A decisão de organizar o desenvolvimento em marcos com documentação contínua se mostrou útil: o diário técnico em \texttt{DEVELOPMENT\_DataAnalyzer.md} registrou as motivações por trás de cada mudança -- inclusive a decisão de não usar \texttt{StatisticsDataFile\_if} -- e facilitou a revisão entre os membros. O comando \texttt{-{}-demo} com 51 testes automatizados garantiu que nenhum marco posterior quebrou os anteriores.

O componente está pronto para ser consumido por outros módulos do simulador via \texttt{TraitsTools<DataAnalyzer\_if>::Implementation}, e serve de base direta para futuros painéis gráficos de análise de resultados.

\newpage

% ---- References ----
\bibliography{references}

\end{document}
